%============================================================
% roadmap.tex --- the whole book at a glance (generated by tools/outline.py)
%============================================================
\chapter{Roadmap of the Book}

\section*{Six Parts, One Staircase}

\begin{diagrambox}
\begin{tikzpicture}[x=2.45cm, y=1.05cm,
    stair/.style={rectangle, rounded corners=2pt, draw=#1, fill=#1!12, line width=1pt,
                  align=center, font=\scriptsize, minimum width=2.3cm, minimum height=1.6cm}]
    \node[stair=successgreen] at (0,0) {\textbf{Part I}\\Foundations\\Ch 1--6\\{\tiny Beginner -- 1700}};
    \node[stair=primaryblue] at (1,1) {\textbf{Part II}\\Core\\Ch 7--12\\{\tiny 1300 -- 2400}};
    \node[stair=softpurple] at (2,2) {\textbf{Part III}\\XOR algebra\\Ch 13--17\\{\tiny 1500 -- 3200}};
    \node[stair=accentorange] at (3,3) {\textbf{Part IV}\\Masks \& words\\Ch 18--22\\{\tiny 1600 -- 2800}};
    \node[stair=warningred] at (4,4) {\textbf{Part V}\\Problem solving\\Ch 23--27\\{\tiny 1500 -- 2600}};
    \node[stair=deepblue] at (5,5) {\textbf{Part VI}\\Advanced\\Ch 28--31\\{\tiny 2200 -- 3500}};
    \draw[flow, line width=1.5pt] (-0.45,-1.1) -- (5.45,4.1);
\end{tikzpicture}
\figcap{Each part assumes fluency with everything below it. Do not skip steps.}
\end{diagrambox}

\section*{Chapter List}

Status: \yes~written \quad $\circ$~planned. For Part~I the \emph{solved} problems of every pattern are listed; every problem appears in exactly \textbf{one} place in the book, and no problem used here is used in any of the seven companion books (array, DP, graph, tree, string, recursion, stack--queue--heap).

\begingroup
\small
\renewcommand{\arraystretch}{1.25}

\subsection*{Part I --- Foundations (Beginner $\to$ CF 1700)}
\begin{tabularx}{\linewidth}{@{}c c >{\raggedright\arraybackslash}p{3.6cm} X@{}}
\toprule
& \textbf{Ch} & \textbf{Title} & \textbf{Patterns (solved problems)} \\
\midrule
$\circ$ & 1 & Binary Representation & Bases and conversions (Convert a Number to Hexadecimal, Convert to Base -2) $\cdot$ Reading the bits of a number (Binary Gap, Binary Number with Alternating Bits) $\cdot$ Fixed width: complement and reversal (Number Complement, Reverse Bits) $\cdot$ Signed integers and overflow (Divide Two Integers, Integer Replacement) $\cdot$ Binary strings as numbers (Binary Prefix Divisible By 5, Concatenation of Consecutive Binary Numbers) \\
$\circ$ & 2 & Bitwise Operators and Their Algebra & Test, set, clear and toggle (Minimum Bit Flips to Convert Number, Minimum Flips to Make a OR b Equal to c) $\cdot$ Shifts (Number of Steps to Reduce a Number to Zero, Number of Even and Odd Bits) $\cdot$ Identities between arithmetic and bitwise operators (Sum of Two Integers, XORwice) $\cdot$ Operators over a range of numbers (Bitwise AND of Numbers Range, And Then There Were K) $\cdot$ Choosing a number under bitwise constraints (Minimize XOR, Boboniu and Bit Operations) \\
$\circ$ & 3 & Classic Bit Tricks & Powers of two (Power of Four, Odd Divisor) $\cdot$ The lowest set bit and carries (Minimum Operations to Reduce an Integer to 0, Johnny and Another Rating Drop) $\cdot$ Highest and lowest bit positions (Smallest Number With All Set Bits, Check if Bitwise OR Has Trailing Zeros) \\
$\circ$ & 4 & Counting Bits and Parity & Popcount (Number of 1 Bits, Sort Integers by The Number of 1 Bits) $\cdot$ Hamming distance (Hamming Distance, Total Hamming Distance) $\cdot$ Set bits over a range of numbers (Counting Bits, Prime Number of Set Bits in Binary Representation) $\cdot$ Counts per bit position (Largest Combination With Bitwise AND Greater Than Zero, And It's Non-Zero) \\
$\circ$ & 5 & XOR Fundamentals & XOR of a range of integers (XOR Operation in an Array, MEXor Mixup) $\cdot$ Undoing XOR (Decode XORed Array, Find The Original Array of Prefix Xor) $\cdot$ XOR equations on neighbours (Neighboring Bitwise XOR, Decode XORed Permutation) $\cdot$ XOR of pairs and triplets (Count Triplets That Can Form Two Arrays of Equal XOR, Bitwise XOR of All Pairings) $\cdot$ Cancellation arguments (AGAGA XOOORRR, Find Xor-Beauty of Array) \\
$\circ$ & 6 & Sets and Subsets as Bitmasks & Enumerating all subsets (Petr and a Combination Lock, Sum of All Subset XOR Totals) $\cdot$ Masks of chosen items (Maximum Rows Covered by Columns, Maximum Length of a Concatenated String with Unique Characters) $\cdot$ Masks of letters and positions (Maximum Product of the Length of Two Palindromic Subsequences, Maximum Score Words Formed by Letters) \\
\bottomrule
\end{tabularx}

\subsection*{Part II --- Core Techniques (CF 1300 $\to$ CF 2400)}
\begin{tabularx}{\linewidth}{@{}c c >{\raggedright\arraybackslash}p{3.6cm} X@{}}
\toprule
& \textbf{Ch} & \textbf{Title} & \textbf{Patterns} \\
\midrule
$\circ$ & 7 & Greedy over Bits: Deciding from the Highest Bit & Maximum AND or XOR of a pair $\cdot$ Answer bit by bit with a check $\cdot$ Minimising under constraints $\cdot$ Lexicographic bit choices \\
$\circ$ & 8 & Per-Bit Independence and Contribution & Sums over pairs $\cdot$ Choosing x bit by bit $\cdot$ Products of bitwise expressions $\cdot$ Independent one-bit subproblems \\
$\circ$ & 9 & AND and OR: Monotonicity, Submasks and Closure & Raising AND with a budget $\cdot$ Submask relations $\cdot$ OR closure $\cdot$ AND chains \\
$\circ$ & 10 & Bitwise Equations and Identities & Sum and XOR given $\cdot$ Equations with carries $\cdot$ Counting solutions $\cdot$ OR and AND given \\
$\circ$ & 11 & Gray Codes, Hypercubes and Bit Orders & Gray-code ranks $\cdot$ Walks on the hypercube $\cdot$ Bit-reversal and other orders $\cdot$ Hamming balls \\
$\circ$ & 12 & Bits in Number Theory & Parity of binomial coefficients $\cdot$ Powers of two and odd parts $\cdot$ Factorials in base 2 $\cdot$ Binary algorithms \\
\bottomrule
\end{tabularx}

\subsection*{Part III --- XOR Algebra (CF 1500 $\to$ CF 3200)}
\begin{tabularx}{\linewidth}{@{}c c >{\raggedright\arraybackslash}p{3.6cm} X@{}}
\toprule
& \textbf{Ch} & \textbf{Title} & \textbf{Patterns} \\
\midrule
$\circ$ & 13 & Prefix XOR and XOR Hashing & Equal prefix XORs $\cdot$ XOR on tree paths $\cdot$ Zobrist hashing $\cdot$ Parity of occurrences \\
$\circ$ & 14 & Linear Basis over GF(2) (XOR Basis) & Maximum and representable XOR $\cdot$ Rank and counting $\cdot$ Bases of prefixes and ranges $\cdot$ Merging bases $\cdot$ k-th smallest XOR \\
$\circ$ & 15 & Gaussian Elimination and Matrices over GF(2) & XOR systems $\cdot$ Toggle puzzles $\cdot$ Counting solutions $\cdot$ Matrix powers over GF(2) \\
$\circ$ & 16 & Walsh--Hadamard and Bitwise Convolutions & XOR convolution $\cdot$ AND and OR convolutions $\cdot$ Transforms in counting problems $\cdot$ Powers of a transform \\
$\circ$ & 17 & Nim and XOR Games & Nim sums $\cdot$ Staircase Nim $\cdot$ Moore's Nim $\cdot$ Grundy values with bit patterns \\
\bottomrule
\end{tabularx}

\subsection*{Part IV --- Bitmask Techniques \& Word Parallelism (CF 1600 $\to$ CF 2800)}
\begin{tabularx}{\linewidth}{@{}c c >{\raggedright\arraybackslash}p{3.6cm} X@{}}
\toprule
& \textbf{Ch} & \textbf{Title} & \textbf{Patterns} \\
\midrule
$\circ$ & 18 & Advanced Mask Enumeration & Submasks $\cdot$ Brute force over masks $\cdot$ k-subsets $\cdot$ Supersets \\
$\circ$ & 19 & Masks as Sets: Compatibility and Intersections & Compatibility $\cdot$ Requirements $\cdot$ Counting pairs by masks $\cdot$ Feature masks \\
$\circ$ & 20 & Bitsets & Reachable sums $\cdot$ Intersections $\cdot$ Transitive closure $\cdot$ Dynamic bitsets \\
$\circ$ & 21 & Word-Level Parallelism and Low-Level Bit Hacks & Parallel counting $\cdot$ Byte tricks $\cdot$ Morton codes $\cdot$ Bit-sliced arithmetic \\
$\circ$ & 22 & Bit Trees: 64-ary Fast Sets and Packed Structures & Successor and predecessor $\cdot$ Packed arrays $\cdot$ Fast sets in sweeps $\cdot$ Bit-tree tricks \\
\bottomrule
\end{tabularx}

\subsection*{Part V --- Bits in Problem Solving (CF 1500 $\to$ CF 2600)}
\begin{tabularx}{\linewidth}{@{}c c >{\raggedright\arraybackslash}p{3.6cm} X@{}}
\toprule
& \textbf{Ch} & \textbf{Title} & \textbf{Patterns} \\
\midrule
$\circ$ & 23 & Randomisation and Hashing with Bits & Robust hashing $\cdot$ Random masks $\cdot$ Hashing states $\cdot$ Probabilistic structures \\
$\circ$ & 24 & Counting over Binary Prefixes & Free suffixes $\cdot$ Bounded XOR $\cdot$ Sums over ranges $\cdot$ Counting pairs under a bound \\
$\circ$ & 25 & Arithmetic Meets Bitwise: Carries and Borrows & No-carry additions $\cdot$ Borrows $\cdot$ Mixed inequalities $\cdot$ Carry chains \\
$\circ$ & 26 & Interactive Bit Problems & Sums from OR and AND $\cdot$ XOR queries $\cdot$ Bit by bit $\cdot$ Budgets and adversaries \\
$\circ$ & 27 & Constructive Bit Problems & Fixing the total XOR $\cdot$ High-bit reservations $\cdot$ Bitwise permutations $\cdot$ Impossibility proofs \\
\bottomrule
\end{tabularx}

\subsection*{Part VI --- Advanced \& Expert Topics (CF 2200 $\to$ CF 3500)}
\begin{tabularx}{\linewidth}{@{}c c >{\raggedright\arraybackslash}p{3.6cm} X@{}}
\toprule
& \textbf{Ch} & \textbf{Title} & \textbf{Patterns} \\
\midrule
$\circ$ & 28 & Expert XOR Algebra & Range bases $\cdot$ Bases on paths $\cdot$ Bases over time $\cdot$ Bases plus counting \\
$\circ$ & 29 & Bit-Level Optimisation and Performance & Built-ins $\cdot$ Branchless code $\cdot$ Packing $\cdot$ Measuring \\
$\circ$ & 30 & Expert Mixed Bit Problems & Mixed set 2400 $\cdot$ Mixed set 2600 $\cdot$ Mixed set 2800 $\cdot$ Mixed set 3000+ \\
$\circ$ & 31 & Contest Mastery & Mixed set 2800 $\cdot$ Mixed set 3000 $\cdot$ Mixed set 3200 $\cdot$ Mixed set 3400+ \\
\bottomrule
\end{tabularx}
\endgroup
